\documentclass[11pt]{article}

% Page layout
\usepackage[letterpaper, margin=1in]{geometry}
\setlength{\headheight}{14pt}
\usepackage{parskip}
\setlength{\parindent}{0pt}

% Typography
\usepackage{fancyvrb}
\usepackage{booktabs}
\usepackage{amsmath}
\usepackage{amssymb}
\usepackage{amsthm}

% Code listings
\usepackage{listings}
\lstset{
    basicstyle=\ttfamily\small,
    columns=fullflexible,
    keepspaces=true,
    showstringspaces=false,
    breaklines=true,
    frame=none,
    xleftmargin=2em,
}

% Hyperlinks
\usepackage[colorlinks=true, linkcolor=blue, urlcolor=blue, citecolor=blue]{hyperref}

% Headers and footers
\usepackage{fancyhdr}
\pagestyle{fancy}
\fancyhf{}
\fancyhead[L]{\textit{Verified Refinement-Based Design Flow}}
\fancyhead[R]{\thepage}
\renewcommand{\headrulewidth}{0.4pt}

% Custom commands (matching csp.tex / cspbuild.tex)
\newcommand{\keyword}[1]{\texttt{\textbf{#1}}}
\newcommand{\nonterminal}[1]{\textit{#1}}

% Theorem environments
\newtheorem{theorem}{Theorem}[section]
\newtheorem{lemma}[theorem]{Lemma}
\newtheorem{definition}[theorem]{Definition}
\newtheorem{proposition}[theorem]{Proposition}

\title{\textbf{Verified Refinement-Based Design Flow \\
for the Fulcrum CSP Toolchain} \\[0.5em]
\large From Specification to Silicon via Compositional Proof}
\author{Mika Nystr\"om \and Claude (Anthropic)}
\date{DRAFT Version 0.1\\
\today}

\begin{document}

\maketitle

\tableofcontents
\newpage

%======================================================================
\section{Introduction}
\label{sec:intro}
%======================================================================

Asynchronous circuit design promises compelling advantages---lower
power, better modularity, natural tolerance of process variation---but
it imposes a correctness burden that synchronous design avoids.
Without a global clock to synchronize all components, the designer
must ensure that every handshake protocol is correct, every pipeline
stage communicates in the right order, and no combination of
interleavings leads to deadlock.  Getting this right is the central
challenge of the field.

The standard approach to this challenge is \emph{post-hoc model
checking}: design the circuit, then feed it to a state-space explorer
that exhaustively checks for deadlock.  FDR4~\cite{fdr4}, the
canonical tool for CSP refinement checking, works this way.  So does
the MDD-based deadlock checker we built for the Fulcrum CSP toolchain,
which scales to over a thousand dining philosophers using symbolic
saturation~\cite{fv-report}.

Post-hoc model checking works, but it is fundamentally the wrong
approach.  Here is why:

\begin{enumerate}
\item \textbf{It reconstructs what the designer already knows.}
  A skilled designer who decomposes a specification into pipeline
  stages does so \emph{because} the decomposition preserves the
  protocol.  The designer's knowledge of \emph{why} the design is
  correct is lost when the design is handed to a model checker, which
  must rediscover the correctness argument by brute-force exploration
  of the state space.

\item \textbf{It does not scale to the design, but to the state space.}
  The cost of model checking depends on the product of all process
  state spaces, not on the complexity of the design itself.
  A clean, well-structured design with 50~processes may have the same
  state-space explosion as a tangled one.  The model checker does not
  exploit the structure that makes the design correct.

\item \textbf{It provides no guidance for repair.}
  When the model checker finds a deadlock, it provides a
  counterexample trace---a sequence of events leading to the deadlocked
  state.  This trace tells the designer \emph{that} something is wrong,
  but not \emph{what} to fix.  The designer must still reason about
  protocols to find the repair, then re-run the model checker to
  verify it.

\item \textbf{It cannot verify implementations.}
  Model checking operates on the CSP-level abstraction.
  It cannot verify that the SystemVerilog or production-rule
  implementation faithfully realizes the CSP specification.
  A separate verification gap exists between the CSP model and the
  circuit.
\end{enumerate}

The thesis of this report is that correctness proof obligations should
be generated \emph{during} the design process, not reconstructed
afterward.  Each design step---decomposing a process, refining a data
type, inserting a pipeline stage, mapping to an implementation
technology---should carry a \emph{local} proof obligation that a tool
can check.  If every step is locally correct, the overall design is
globally correct by composition.

This is not a new idea.  It is the essence of Alain Martin's
synthesis methodology~\cite{martin1990}, which derives circuits from
specifications via a sequence of semantics-preserving
transformations.  It is the principle behind Event-B's refinement
calculus~\cite{abrial2010}.  It underlies session types for
concurrent programming~\cite{honda1998,kobayashi2006}.  What is new
here is the application of this principle to a \emph{concrete,
existing} toolchain---the Fulcrum CSP tools---with a concrete proposal
for the annotations, checks, and tools needed to make it work.

The scope of this report covers the entire path from initial
specification to final implementation:
\begin{itemize}
\item Bundled-data (BD) asynchronous circuits via \texttt{csp2verilog}
\item Synchronous RTL via dehandshaking
\item Quasi-delay-insensitive (QDI) pipeline circuits via
  \texttt{csp2gma} and production rules
\end{itemize}
We describe what annotations the designer adds to \texttt{.csp} and
\texttt{.sys} files, what the tool checks at each step, what tools
exist, and what needs to be built.

%======================================================================
\section{The Current Toolchain}
\label{sec:current}
%======================================================================

%----------------------------------------------------------------------
\subsection{Fulcrum CSP Language}
%----------------------------------------------------------------------

The Fulcrum CSP dialect (sometimes called CAST) is a hardware-oriented
process description language in the tradition of Hoare's
CSP~\cite{hoare1985} and Martin's CHP~\cite{martin1990}; see the
language reference manual~\cite{nystrom2025csp} for the complete
syntax and semantics.  A process
is a sequential program that communicates with its environment via
directional channels:

\begin{Verbatim}[samepage=true]
int(32) x;
*[ L?x ;                        // receive from channel L
   x = x + 1 ;
   R!x                          // send on channel R
]
\end{Verbatim}

Key language features include:
\begin{description}
\item[Fixed-width data.]
  Types are \texttt{int(\nonterminal{W})} for $W$-bit integers,
  \texttt{bool}, arrays, and structures.  All widths are statically
  known, matching the hardware target.

\item[Directional channels.]
  Channel operations are \texttt{L?x} (receive) and \texttt{R!x}
  (send).  Each channel has a statically declared width and is used by
  exactly two processes---one sender, one receiver.

\item[Deterministic selection.]
  The guarded command \texttt{[\,]} selects among guards based on
  channel readiness (\emph{probes}) or data conditions:
\begin{Verbatim}[samepage=true]
*[[ #L -> L?x ; R!x
  :  #S -> S?y ; T!y
]]
\end{Verbatim}

\item[Nondeterministic selection.]
  The \texttt{:[\,]:} construct provides internal (nondeterministic)
  choice, used when the environment cannot observe which branch is
  taken.

\item[Repetition.]
  The \texttt{*[\,]} construct is an infinite loop---the normal
  operating mode of a hardware process.

\item[Structures.]
  The \keyword{structure} declaration packs multiple fields for
  transmission over a single channel:
\begin{Verbatim}[samepage=true]
structure EmData = (int(32) result; int(5) dest; int(1) rw);
EmData em;
em::result = alu_out;
em_data!em
\end{Verbatim}
\end{description}

%----------------------------------------------------------------------
\subsection{The \texttt{.sys} System Description Format}
%----------------------------------------------------------------------

Multi-process systems are described in \texttt{.sys} files, which
declare the channels, process types, and instances.  Each process type
references an external \texttt{.csp} file and declares its ports with
directions and widths:

\begin{Verbatim}[samepage=true]
system ProdCons;
  var c : channel(32);

  process Producer = "producer.csp"
    port out R : channel(32);

  process Consumer = "consumer.csp"
    port in L : channel(32);
begin
  var p : Producer(R => c);
  var q : Consumer(L => c);
end.
\end{Verbatim}

The instance section binds each process port to a system-level
channel using \texttt{\nonterminal{port} => \nonterminal{channel}}
syntax.  The system builder validates that all widths match, all ports
are bound, and each channel has exactly one sender and one receiver.

%----------------------------------------------------------------------
\subsection{Existing Tools}
%----------------------------------------------------------------------

The Fulcrum CSP toolchain currently comprises the following tools:

\begin{description}
\item[\texttt{cspfe} (parser).]
  Lex/Yacc-based parser that reads \texttt{.csp} source files and
  emits S-expression ASTs.  The AST is defined by the
  \texttt{cspgrammar} Modula-3 library.

\item[\texttt{cspc} (compiler/driver).]
  The main tool.  With \texttt{-scm} it enters a Scheme REPL (via the
  embedded \texttt{mscheme} interpreter) that drives a 9-pass
  compilation pipeline: parsing, type checking, optimization, and code
  generation of Modula-3 state machines.

\item[\texttt{cspbuild} (system builder).]
  A Scheme-level build system that reads \texttt{.sys} files,
  orchestrates compilation of each process, and generates simulator
  binaries.  Entry point: \texttt{(build-system!\ "file.sys")}.

\item[\texttt{csp2verilog} (BD synthesis).]
  Java tool that converts CSP specifications to synthesizable
  SystemVerilog with bundled-data handshake protocols.  Uses a visitor
  pattern over the CAST AST.  Supports coverage probes and state
  covergroups.

\item[\texttt{Cast2Verilog} (full design synthesis).]
  Comprehensive Java tool that handles both structural (PRS-based) and
  behavioral (CSP-based) CAST cell designs.  Generates hierarchical
  Verilog, testbenches, and co-simulation interfaces.

\item[\texttt{csp2gma} (QDI path).]
  Java tool that converts CSP to GMA (Graph-based Modeling and
  Analysis) format, an intermediate representation for
  quasi-delay-insensitive circuit synthesis.

\item[\texttt{slacker} (slack optimization).]
  Java tool that uses linear programming (via GLPK) to compute optimal
  buffer placements and slack assignments.  Formulates the
  slack-matching problem as an LP with per-channel constraints and a
  cost function minimizing area overhead.

\item[Formal verification.]
  The MDD-based deadlock checker (\texttt{check-deadlock-saturation!})
  uses Multi-valued Decision Diagrams with a saturation algorithm for
  symbolic state-space exploration.  It scales to 1024-process dining
  philosopher systems.  An explicit-state checker with partial-order
  reduction (\texttt{check-deadlock!}) handles smaller systems with
  counterexample generation.
\end{description}

%----------------------------------------------------------------------
\subsection{The MiniMIPS as Running Example}
\label{sec:minimips}
%----------------------------------------------------------------------

Throughout this report we use the MiniMIPS sequential processor as a
running example.  It is a 5-stage pipeline (Fetch, Decode, Execute,
MemStage, Writeback) with two server processes (RegFile, DataMem)---7
processes connected by 26 channels.  The system topology is declared
in \texttt{minimips.sys}:

\begin{Verbatim}[samepage=true]
system MiniMIPS;
  // Pipeline channels: Fetch -> Decode
  var fd_pc    : channel(32);
  var fd_instr : channel(32);

  // Pipeline channels: Decode -> Execute
  var de_pc     : channel(32);
  var de_instr  : channel(32);
  var de_rs_val : channel(32);
  var de_rt_val : channel(32);
  var de_imm    : channel(32);

  // Execute -> MemStage (struct-packed)
  var em_data : channel(107);

  // MemStage -> Writeback
  var mw_result : channel(32);
  var mw_dest   : channel(5);
  // ... (20 more channels)

  // Feedback: Writeback -> Fetch
  var wb_taken  : channel(1);
  var wb_target : channel(32);
  // ...
begin
  var fetch     : Fetch    (fd_pc => fd_pc, ...);
  var decode    : Decode   (...);
  var execute   : Execute  (...);
  var memstage  : MemStage (...);
  var writeback : Writeback(...);
  var regfile   : RegFile  (...);
  var datamem   : DataMem  (...);
end.
\end{Verbatim}

The MiniMIPS exhibits the key structural features that our verification
framework must handle: linear pipeline flow, feedback channels (branch
resolution), server processes with multiple clients (RegFile serves
both Decode ports), and structure-packed wide channels.

%----------------------------------------------------------------------
\subsection{Current Tool Flow}
%----------------------------------------------------------------------

The current path from specification to implementation is:

\begin{Verbatim}[samepage=true]
                  .csp sources        .sys topology
                       |                    |
                       v                    v
                    cspfe              sys parser
                       |                    |
                       v                    v
                  S-expr AST          system graph
                       |                    |
              +--------+--------+     +-----+-----+
              |        |        |     |           |
              v        v        v     v           v
           cspc    csp2verilog  csp2gma  slacker   MDD checker
              |        |        |     |           |
              v        v        v     v           v
        M3 simulator  SV BD    GMA   slack LP   deadlock
        binary        netlist  IR    solution   result
\end{Verbatim}

The critical observation is that verification is a \emph{separate
branch}---it operates on the same source but shares no information
with the synthesis path.  The model checker does not know that the
designer structured the pipeline to preserve protocols; the synthesis
tool does not know that the model checker verified deadlock freedom.
The two paths are disconnected.

%======================================================================
\section{The Verified Refinement Framework}
\label{sec:framework}
%======================================================================

This section presents the theoretical framework for verified
refinement-based design.  The core idea is that each channel in a
system carries a \emph{protocol type} that specifies the sequence of
communications permitted on that channel.  If every process conforms
to its channel protocols, and the protocols at each channel's two
endpoints match, then the composed system is deadlock-free
\emph{by construction}.  Design proceeds by a sequence of
refinement steps, each of which preserves protocol conformance and is
checked locally.

%----------------------------------------------------------------------
\subsection{Protocol Types on Channels}
\label{sec:protocol-types}
%----------------------------------------------------------------------

A \emph{protocol type} describes the permitted sequence of
values communicated on a channel.  Since FSCP channels are
\emph{unidirectional}---each channel has exactly one sender and one
receiver---a protocol on a single channel describes only the sequence
of values transferred, not a mix of sends and receives.  Multi-step
interactions such as request-response are implemented with multiple
channels (one per direction), and their coordination is captured by
the channel dependency graph (Section~\ref{sec:deadlock-freedom}),
not by the protocol type of any single channel.

\begin{definition}[Protocol Type]
\label{def:protocol}
A protocol type $\pi$ is defined by the grammar:
\[
\pi ::= \tau \mid
        \pi_1 \mathbin{;} \pi_2 \mid
        \pi_1 \mid \pi_2 \mid
        \pi^{*} \mid \texttt{end}
\]
where $\tau$ is a data type (possibly with value constraints),
$\mathbin{;}$ is sequencing, $\mid$ is choice, ${}^{*}$ is
repetition, and \texttt{end} is the terminated protocol (no further
communication).  Each element $\tau$ in the protocol represents one
handshake---one value transferred from sender to receiver.
\end{definition}

We propose the following concrete syntax for protocol annotations in
\texttt{.sys} files:

\begin{Verbatim}[samepage=true]
var c : channel(32) protocol { v }*;
\end{Verbatim}

This declares a channel that repeatedly transfers a 32-bit value---one
value per handshake cycle.  The protocol is the same from both
endpoints' perspective (the sender sends $v$, the receiver receives
$v$; the direction is already fixed by the channel's
\keyword{port}~\keyword{in}/\keyword{out} declaration).

More examples:

\begin{Verbatim}[samepage=true]
// Simple data channel: one value per handshake
var data : channel(32) protocol { v }*;

// Alternating values with a constraint
var tagged : channel(33) protocol { v : v[32]=0 ; v : v[32]=1 }*;

// Finite protocol: exactly N values, then done
var init : channel(32) protocol { v }^N ; end;
\end{Verbatim}

For most channels in typical FSCP designs, the protocol is simply
\texttt{\{ v \}*}---an infinite stream of values, one per handshake.
Richer protocols arise when there are value constraints (e.g., a
tag bit that alternates) or when the channel carries a finite
sequence (e.g., an initialization channel).

\subsubsection{Multi-channel protocols and dependency}

Request-response patterns---such as the RegFile's read port---use
\emph{two} channels: one for the request, one for the response.
Each channel individually has the simple protocol \texttt{\{ v \}*}.
The constraint that a request must precede its response is not a
property of any single channel's protocol; it is a \emph{dependency}
between channels, captured in the channel dependency graph
(Definition~\ref{def:cdg}).  This clean separation is a direct
consequence of FSCP's unidirectional channel model.

\subsubsection{Protocol matching}

Since both endpoints of a channel see the same sequence of values (the
sender produces them, the receiver consumes them), the protocol is
shared, not dualized.  The tool checks that both the sender and the
receiver agree on the channel's protocol type---that is, the sequence
of values the sender produces must match the sequence the receiver
expects.  This replaces the duality condition found in session type
systems, which is needed only for bidirectional channels.

\subsubsection{Protocol conformance}

A process \emph{conforms} to its port protocols if, for every
reachable state, the next communication on each port is consistent
with the protocol type declared for that port.  Formally:

\begin{definition}[Protocol Conformance]
\label{def:conformance}
Process $P$ with ports $p_1, \ldots, p_n$ having protocol types
$\pi_1, \ldots, \pi_n$ \emph{conforms} to its protocols if, for every
reachable state $s$ of $P$ and every port $p_i$, the sequence of
communications performed by $P$ on $p_i$ up to state $s$ is a prefix
of some word in $\mathcal{L}(\pi_i)$, the language of protocol~$\pi_i$.
\end{definition}

Protocol conformance is a \emph{local} property---it can be checked
for each process independently, without composing it with the rest of
the system.  This is the key to scalability.

\subsubsection{Relationship to session types}

Protocol types as defined here are a radical simplification of
\emph{binary session types}~\cite{honda1998}.  Session types assign
dual types to the two endpoints of a bidirectional channel ($!\tau$
for send, $?\tau$ for receive).  Our protocol types describe only
the value sequence on a unidirectional channel; the direction is
fixed by the port declaration, not encoded in the type.

\begin{center}
\begin{tabular}{ll}
\toprule
\textbf{Session types} & \textbf{FSCP protocol types} \\
\midrule
Output type $!\tau.\sigma$ / Input type $?\tau.\sigma$
  & Value transfer $\tau \mathbin{;} \pi$ \\
  & (direction fixed by port decl.) \\
Selection $\oplus\{l_i: \sigma_i\}$ & Choice $\pi_1 \mid \pi_2$ \\
Branching $\&\{l_i: \sigma_i\}$ & Choice (same; no duality) \\
Recursion $\mu X.\sigma$ & Repetition $\pi^{*}$ \\
End & \texttt{end} \\
Duality $!\tau \leftrightarrow ?\tau$ & Not needed (unidirectional) \\
Delegation (channel passing) & Not needed (static topology) \\
\bottomrule
\end{tabular}
\end{center}

The key simplifications relative to full session types are:
(1)~channels are unidirectional, eliminating duality from the type
system entirely; (2)~channels are point-to-point with fixed-width
data, eliminating delegation and polymorphism; (3)~multi-channel
coordination (e.g., request before response) is handled by the
channel dependency graph rather than by nesting send/receive in a
single type.  These simplifications make the checking problem
decidable by finite automaton inclusion rather than the more complex
type-checking algorithms needed for general session types.

%----------------------------------------------------------------------
\subsection{Compositional Deadlock-Freedom}
\label{sec:deadlock-freedom}
%----------------------------------------------------------------------

The protocol type system enables a compositional deadlock-freedom
theorem that avoids global state-space exploration entirely.

\begin{definition}[Channel Dependency Graph]
\label{def:cdg}
The \emph{channel dependency graph} $G = (V, E)$ of a system has one
vertex per channel.  There is an edge $(c_1, c_2)$ if some process
performs an operation on $c_2$ that depends on the result of an
operation on $c_1$---that is, the process must complete a communication
on $c_1$ before it can initiate a communication on $c_2$.
\end{definition}

\begin{theorem}[Compositional Deadlock Freedom]
\label{thm:deadlock-free}
Let $S = P_1 \| P_2 \| \cdots \| P_n$ be a system of processes
communicating on channels $c_1, \ldots, c_m$.  If:
\begin{enumerate}
\item[(a)] Every process $P_i$ conforms to its port protocols
  (Definition~\ref{def:conformance}).
\item[(b)] Every channel $c_j$ has matching protocols at its two
  endpoints (the sender's value sequence matches what the receiver
  expects).
\item[(c)] The channel dependency graph $G$ (Definition~\ref{def:cdg})
  is acyclic.
\end{enumerate}
Then $S$ is deadlock-free.
\end{theorem}

\begin{proof}[Proof sketch]
Suppose for contradiction that $S$ reaches a deadlocked state~$s$
where every process is blocked waiting on some channel.  Since $G$ is
acyclic, it has a topological ordering $c_{i_1}, \ldots, c_{i_m}$.
Consider the channel $c_{i_1}$ that is first in this ordering.
The process $P$ that is blocked waiting on $c_{i_1}$ has no prior
dependency (by the topological ordering), so it is ready to
communicate on $c_{i_1}$.  By protocol conformance~(a), $P$'s next
action on $c_{i_1}$ matches the protocol.  By protocol matching~(b),
the other endpoint's process $Q$ expects the same value sequence.
By protocol conformance of $Q$, $Q$ is ready for this action (since
$Q$'s prior
dependencies are all on channels later in the ordering, which
contradicts $c_{i_1}$ being first).  But then $P$ and $Q$ can
communicate on $c_{i_1}$, contradicting the assumption that $s$ is
deadlocked. \qed
\end{proof}

\subsubsection{Handling cycles}

Many real systems have cyclic channel dependencies.  The MiniMIPS has
a cycle through the Writeback$\to$Fetch feedback channels.  When the
channel dependency graph contains cycles, the compositional theorem
does not apply directly, and the designer must break the cycle.
There are several strategies:

\begin{description}
\item[Slack insertion.]
  Adding a buffer (slack) on a channel in the cycle can break the
  dependency by ensuring that one endpoint is always ready.  A channel
  with slack $d \geq 1$ allows the sender to proceed without waiting
  for the receiver (up to $d$ pending values), breaking the
  synchronous dependency.

\item[Asymmetric ordering.]
  The classic solution for dining philosophers: one process in the
  cycle acquires its channels in a different order, breaking the
  circular wait.

\item[Localized model checking.]
  For cycles that cannot be broken structurally, the system can
  extract the cycle subgraph and apply model checking only to the
  processes involved in the cycle.  This is still compositional: the
  acyclic portion is verified by the protocol type system, and only the
  cyclic core requires state-space exploration.
\end{description}

\subsubsection{Relationship to prior work}

The compositional deadlock-freedom theorem draws on several lines of
work:

\begin{itemize}
\item Kobayashi's type systems for deadlock-freedom in the
  $\pi$-calculus~\cite{kobayashi2006} assign ``obligation levels'' and
  ``capability levels'' to channels to ensure acyclic dependency.  Our
  channel dependency graph is the explicit version of this implicit
  ordering.

\item Padovani's deadlock-freedom by typing~\cite{padovani2014}
  extends Kobayashi's work to more general topologies.  Our approach
  is less general (we require explicit acyclicity checking) but more
  directly applicable to hardware designs where the topology is
  statically known.

\item The ``communication closure'' condition in rely-guarantee
  reasoning~\cite{jones1983} ensures that processes can be verified
  independently.  Protocol conformance is our version of this
  condition.
\end{itemize}

%----------------------------------------------------------------------
\subsection{Refinement Steps}
\label{sec:refinement-steps}
%----------------------------------------------------------------------

Design proceeds by a sequence of refinement steps, each of which
transforms the system while preserving its essential properties.
Each step generates a local proof obligation that a tool can check.

\subsubsection{Process decomposition}

The most common refinement step is splitting a single process into
two or more parallel sub-processes connected by new internal channels.

\begin{definition}[Process Decomposition]
A process $P$ with ports $\{p_1, \ldots, p_n\}$ is decomposed into
$P_1 \| P_2$ with:
\begin{itemize}
\item $P_1$ has ports $\{p_1, \ldots, p_k\} \cup \{q_1, \ldots, q_r\}$
\item $P_2$ has ports $\{p_{k+1}, \ldots, p_n\} \cup \{\bar{q}_1,
  \ldots, \bar{q}_r\}$
\item $q_1, \ldots, q_r$ are new internal channels connecting $P_1$
  to $P_2$, with $\bar{q}_i$ being the other endpoint of $q_i$.
\end{itemize}
\end{definition}

The proof obligation for process decomposition is \emph{trace
equivalence}: the externally visible behavior of $P_1 \| P_2$ (hiding
the internal channels $q_1, \ldots, q_r$) must equal the behavior of
the original process~$P$.  That is:
\[
P =_{\mathcal{T}} (P_1 \| P_2) \setminus \{q_1, \ldots, q_r\}
\]

In practice, we check a weaker condition that is sufficient for
correctness: $P_1$ and $P_2$ each conform to their port protocols
(including the new internal channels), the internal channels have
matching protocols at both endpoints, and the augmented channel
dependency graph remains acyclic
(or any new cycles are handled by slack or localized checking).

\textbf{Example.} Consider a monolithic pipeline process that reads
from input channel~$L$, processes the data, and writes to output
channel~$R$:

\begin{Verbatim}[samepage=true]
// Monolithic version
int(32) x;
*[ L?x ; x = transform(x) ; R!x ]
\end{Verbatim}

Decompose into two stages:

\begin{Verbatim}[samepage=true]
// Stage 1
int(32) x;
*[ L?x ; x = partial_transform(x) ; M!x ]

// Stage 2
int(32) x;
*[ M?x ; x = remaining_transform(x) ; R!x ]
\end{Verbatim}

The proof obligation is that \texttt{partial\_transform} composed with
\texttt{remaining\_transform} equals \texttt{transform} for all
inputs, and that the internal channel~$M$ has matching protocols
at its two endpoints.

\subsubsection{Data refinement}

Data refinement replaces an abstract data representation with a
concrete one.  In hardware terms, this is the step from ``this channel
carries an integer'' to ``this channel carries a 32-bit two's
complement value.''

\begin{definition}[Data Refinement]
A data refinement replaces type $\tau_{\text{abs}}$ with
$\tau_{\text{conc}}$ where there exists a representation function
$\rho : \tau_{\text{conc}} \to \tau_{\text{abs}}$ that is surjective
(every abstract value has at least one concrete representation).
\end{definition}

The proof obligation is that the protocol is preserved: if the
abstract process sends values $v_1, v_2, \ldots$ satisfying some
data constraint, then the concrete process sends values $w_1, w_2,
\ldots$ such that $\rho(w_i) = v_i$ for all $i$, and the constraint
is still satisfied under $\rho$.

In the Fulcrum CSP context, the primary data refinement is narrowing
from \texttt{int} (simulation type, arbitrary precision) to
\texttt{int(\nonterminal{W})} (synthesis type, $W$~bits).  The tool
must verify that no computation overflows the reduced width, or that
overflow behavior matches the intended semantics.

\subsubsection{Channel splitting}

Channel splitting replaces a single wide channel with multiple
narrower channels.  This is common when a structure-packed channel is
decomposed into its constituent fields.

\begin{Verbatim}[samepage=true]
// Before: single 107-bit structure channel
var em_data : channel(107) protocol {em }*;

// After: three separate channels
var em_result : channel(32) protocol {result }*;
var em_dest   : channel(5)  protocol {dest   }*;
var em_rw     : channel(1)  protocol {rw     }*;
\end{Verbatim}

The proof obligation is protocol equivalence: the sequence of
communications on the split channels, interleaved in the declared
order, must correspond to the original protocol on the packed channel.

\subsubsection{Pipeline insertion (Lines-style)}

Pipeline insertion adds buffer stages to increase throughput.
In Lines' pipelining theory~\cite{lines1998}, a pipeline stage is a
process that receives on one channel and sends on another, with a
specific handshake template (WCHB, PCHB, PCFB, etc.).

\begin{Verbatim}[samepage=true]
// Before: direct connection
var c : channel(32) protocol {v }*;

// After: c split into c1 -> buffer -> c2
var c1 : channel(32) protocol {v }*;
var c2 : channel(32) protocol {v }*;
process Buffer = begin int(32) x; *[ c1?x ; c2!x ] end
    port in c1 : channel(32)
    port out c2 : channel(32);
\end{Verbatim}

The proof obligation is that the buffer preserves the protocol: the
sequence of values on $c_2$ equals the sequence of values on $c_1$
(possibly delayed by the buffer depth).  This is trivially true for
the simple buffer above; for more complex pipeline templates with
handshake reshuffling, the proof involves showing that the template's
internal protocol matches the declared channel protocols.

%----------------------------------------------------------------------
\subsection{Separation of Optimization from Correctness}
\label{sec:separation}
%----------------------------------------------------------------------

A key design principle of the framework is that \emph{optimizations
are verified transformations}, not ad hoc modifications.

\subsubsection{Slack insertion as a verified transformation}

Adding buffers to a channel (increasing its slack from $d$ to $d+k$)
is a transformation that preserves protocol conformance.  A buffer
process with slack~$k$ on a channel with protocol~$\pi$ simply relays
values according to $\pi$ with up to $k$ values in flight.  The proof
obligation is that the buffer conforms to~$\pi$ on both its input
and output ports---which is trivially the case for a FIFO buffer.

The \texttt{slacker} tool already computes optimal slack assignments
via linear programming.  In the verified framework, each slack
assignment is a pipeline insertion step with a trivial proof
obligation.

\subsubsection{Slack elasticity}

A system is \emph{slack elastic} if its functional behavior (the
sequence of values produced on each output channel) is independent
of the amount of buffering on its internal
channels~\cite{manohar2001,cardi2011}.  Slack elasticity is a
critical property for the verified refinement framework because it
means that slack insertion and removal are \emph{semantics-preserving
transformations}---the designer (or the \texttt{slacker} tool) can
adjust buffering for performance without affecting correctness.

Not all systems are slack elastic.  A system with data-dependent
control flow that observes buffer occupancy (e.g., through probes
that distinguish empty from non-empty channels) may change behavior
when slack changes.  The sufficient condition for slack elasticity
is that no process uses the \emph{absence} of data on a channel
as a control signal---that is, guards depend only on the
\emph{presence} of data (\texttt{\#L} probes that test whether a
channel has data ready) or on data values, never on whether a
channel is empty.

In the protocol type framework, slack elasticity can be verified as
follows:
\begin{enumerate}
\item For each channel, check that no process guards on the
  \emph{negation} of a probe on that channel (i.e., no
  ``\texttt{\~{}\#L}'' guards that would distinguish slack~0 from
  slack~$>0$).
\item Verify that the channel dependency graph is preserved under
  slack changes (adding buffers does not introduce new dependencies).
\end{enumerate}
If these conditions hold, the system is slack elastic and any slack
assignment computed by \texttt{slacker} preserves the functional
behavior verified at the CSP level.  This provides a clean separation:
the protocol type system verifies functional correctness, slack
elasticity guarantees that performance optimization does not
invalidate the proof, and \texttt{slacker} optimizes throughput within
the verified-correct design.

\subsubsection{Action reordering}

Independent actions within a process may be reordered for optimization
(e.g., to improve pipeline throughput or reduce critical path length).
Two actions are independent if they operate on different channels and
neither reads data produced by the other.

The proof obligation for action reordering is that the reordered
sequence is a valid interleaving of the original sequence with respect
to the channel protocols.  This connects directly to partial-order
reduction (POR) theory~\cite{peled1993}: the independence relation
used for action reordering is the same one used to prune the
state-space search in model checking.

\subsubsection{Technology mapping}

Technology mapping---the step from CSP to gates, production rules, or
RTL---operates \emph{below} the CSP abstraction level.  The CSP-level
verification guarantees protocol correctness; technology mapping must
additionally satisfy timing constraints (setup/hold for synchronous,
delay matching for BD, isochronic fork for QDI) that are outside the
scope of CSP-level reasoning.

The framework treats technology mapping as a separate concern:
the CSP-level proof establishes \emph{functional} correctness (the
right values appear on the right channels in the right order), and
implementation-level tools (static timing analysis, production rule
verification) establish \emph{timing} correctness.

%======================================================================
\section{The Design Flow}
\label{sec:flow}
%======================================================================

This section describes the concrete design flow from specification to
implementation, organized into four stages.  Each stage is illustrated
with examples from the MiniMIPS and simpler designs.

%----------------------------------------------------------------------
\subsection{Stage 1: Specification}
\label{sec:stage1}
%----------------------------------------------------------------------

The design flow begins with an executable CSP specification---either a
Martin-style process description or an external I/O contract.

\subsubsection{Writing the specification}

The specification is an ordinary \texttt{.csp} file that captures the
intended behavior.  It need not be implementable directly; it serves
as the reference against which all subsequent refinements are checked.

\textbf{Example: FIFO buffer specification.}

\begin{Verbatim}[samepage=true]
// fifo_spec.csp — behavioral spec for a 4-deep FIFO
int(32) buf[4];
int(3)  head, tail, count;

*[[ count < 4 && #L -> L?buf[tail] ; tail = (tail+1) % 4 ; count = count+1
  : count > 0 && #R -> R!buf[head] ; head = (head+1) % 4 ; count = count-1
]]
\end{Verbatim}

\textbf{Example: arbiter specification.}

\begin{Verbatim}[samepage=true]
// arbiter_spec.csp — 2-input arbiter
int(32) x;
*[[ #L0 -> L0?x ; R!x
  : #L1 -> L1?x ; R!x
]]
\end{Verbatim}

\subsubsection{Annotating channel protocols}

The designer annotates each channel in the \texttt{.sys} file with its
protocol type.  For the producer-consumer example:

\begin{Verbatim}[samepage=true]
system ProdCons;
  var c : channel(32) protocol {v }*;

  process Producer = "producer.csp"
    port out R : channel(32) conforms { v }*;

  process Consumer = "consumer.csp"
    port in L : channel(32) conforms { v }*;
begin
  var p : Producer(R => c);
  var q : Consumer(L => c);
end.
\end{Verbatim}

The \keyword{protocol} clause on the channel declaration specifies the
protocol from the sender's perspective.  The \keyword{conforms} clause
on each port declaration specifies the protocol that the process
guarantees on that port.  The tool checks:
\begin{enumerate}
\item Each port's \keyword{conforms} protocol matches the channel's
  \keyword{protocol} (the same value sequence, regardless of direction).
\item Each process's code actually conforms to its declared port
  protocols.
\end{enumerate}

For simple data channels (one send per handshake cycle), the protocol
is \texttt{\{ v \}*} and the annotation is straightforward.  For
request-response protocols like the RegFile's read port:

\begin{Verbatim}[samepage=true]
// RegFile read protocol: requester sends address, gets data back
var rs_req  : channel(5)  protocol {addr }*;
var rs_resp : channel(32) protocol {data }*;
\end{Verbatim}

Here the two channels together implement a request-response protocol.
The tool verifies that the requester (Decode) always sends on
\texttt{rs\_req} before receiving on \texttt{rs\_resp}, and that the
responder (RegFile) always receives on \texttt{rs\_req} before sending
on \texttt{rs\_resp}.  This ordering constraint is captured in the
channel dependency graph rather than in the individual protocol types.

\subsubsection{Initial verification}

At this stage, the tool performs three checks:
\begin{enumerate}
\item \textbf{Protocol conformance.} Each process's code conforms to
  its declared port protocols.
\item \textbf{Protocol matching.} Each channel's sender and receiver
  agree on the protocol (same value sequence).
\item \textbf{Dependency acyclicity.} The channel dependency graph is
  acyclic (or cycles are annotated as requiring special treatment).
\end{enumerate}

If all three checks pass, the system is deadlock-free by
Theorem~\ref{thm:deadlock-free}.  If the dependency graph has cycles,
the tool reports them and the designer must either break them
(slack, asymmetry) or mark them for localized model checking.

%----------------------------------------------------------------------
\subsection{Stage 2: Process Decomposition}
\label{sec:stage2}
%----------------------------------------------------------------------

The designer decomposes processes to expose pipeline parallelism, separate
concerns, or match a target micro-architecture.

\subsubsection{The decomposition step}

Suppose the designer starts with a monolithic specification:

\begin{Verbatim}[samepage=true]
// mips_spec.csp — monolithic MIPS: fetch, decode, execute in one process
int(32) pc, instr, rs_val, rt_val, result;
*[
  instr = imem[pc >> 2] ;
  // decode
  rs_val = regs[(instr >> 21) & 0x1F] ;
  rt_val = regs[(instr >> 16) & 0x1F] ;
  // execute
  result = alu(rs_val, rt_val, instr) ;
  // writeback
  regs[(instr >> 11) & 0x1F] = result ;
  pc = pc + 4
]
\end{Verbatim}

The designer decomposes this into the 5-stage pipeline reflected in
\texttt{minimips.sys}.  Each decomposition step introduces new
internal channels.  The \texttt{.sys} file records the decomposition:

\begin{Verbatim}[samepage=true]
system MiniMIPS;
  // New internal channels (from decomposition)
  var fd_pc    : channel(32) protocol {v }*;
  var fd_instr : channel(32) protocol {v }*;
  var de_pc    : channel(32) protocol {v }*;
  // ...

  process Fetch = "fetch.csp"
    port out fd_pc    : channel(32) conforms { v }*
    port out fd_instr : channel(32) conforms { v }*
    port in  wb_taken : channel(1)  conforms { v }*
    port in  wb_target: channel(32) conforms { v }*;

  process Decode = "decode.csp"
    port in  fd_pc     : channel(32) conforms { v }*
    port in  fd_instr  : channel(32) conforms { v }*
    port out de_pc     : channel(32) conforms { v }*
    // ...
\end{Verbatim}

\subsubsection{Proof obligations}

For each decomposition step, the tool checks:

\begin{enumerate}
\item \textbf{Protocol conformance of sub-processes.}
  Each new process (\texttt{fetch.csp}, \texttt{decode.csp}, etc.)
  conforms to its declared port protocols.  This is a local check on
  each process individually.

\item \textbf{Matching protocols on new internal channels.}
  The \texttt{fd\_pc} channel has protocol \texttt{\{ v \}*}; both
  the Fetch end (sender) and the Decode end (receiver) agree on this
  value sequence.

\item \textbf{Dependency acyclicity.}
  The tool constructs the channel dependency graph for the decomposed
  system and checks for cycles.  The MiniMIPS has a cycle:
  \[
  \text{Fetch} \xrightarrow{\text{fd\_*}} \text{Decode}
  \xrightarrow{\text{de\_*}} \text{Execute}
  \xrightarrow{\text{em\_*}} \text{MemStage}
  \xrightarrow{\text{mw\_*}} \text{Writeback}
  \xrightarrow{\text{wb\_*}} \text{Fetch}
  \]
  The designer must annotate this cycle, indicating that it is broken
  by the initial token in the Fetch process (the initial PC value).
\end{enumerate}

\subsubsection{Tool support for decomposition}

The decomposition itself is manual---the designer writes the sub-process
code.  The tool assists by:
\begin{itemize}
\item Automatically inferring the channel dependency graph from the
  process code.
\item Checking protocol conformance.
\item Reporting cycles and suggesting where slack insertion or
  asymmetry could break them.
\item Optionally, running localized model checking on cyclic
  subgraphs.
\end{itemize}

%----------------------------------------------------------------------
\subsection{Stage 3: Data Refinement}
\label{sec:stage3}
%----------------------------------------------------------------------

Data refinement narrows the abstract types used in the specification
to the concrete bit widths required for synthesis.  The Fulcrum CSP
compiler (\texttt{cspc}) already performs substantial data refinement
automatically during its 9-pass compilation pipeline; the verified
framework aims to make the correctness of this refinement explicit
and checkable.

\subsubsection{Existing infrastructure in \texttt{cspc}}

The \texttt{cspc} compiler represents integer types as 5-element
S-expressions:
\begin{Verbatim}[samepage=true]
(integer <const?> <signed?> <bitwidth> <range>)
\end{Verbatim}
where \texttt{bitwidth} is either a BigInt (for \texttt{int(N)}) or
the empty list (for unbounded \texttt{int}).  The three integer type
categories are:
\begin{description}
\item[\texttt{int}] --- infinite-precision signed integer (from $\mathbb{Z}$).
\item[\texttt{int(N)}] --- $N$-bit unsigned, range $0 \ldots 2^N-1$.
\item[\texttt{sint(N)}] --- $N$-bit signed two's complement,
  range $-2^{N-1} \ldots 2^{N-1}-1$.
\end{description}

Width determination is deferred to \textbf{compile6!}, the sixth
compiler pass, which performs:
\begin{enumerate}
\item \textbf{Table construction} (\texttt{make-the-tables}).  Builds
  per-variable tables of all assignments, uses, declared ranges, and
  port (channel) associations.

\item \textbf{Range seeding} (\texttt{seed-integer-ranges!}).
  Initializes the range table from user-declared types.  Channel
  receive operations seed ranges from the channel's bit width.

\item \textbf{Fixed-point interval arithmetic}
  (\texttt{close-integer-ranges!}).  Iterates up to 10 rounds,
  refining each variable's range by computing the union of its
  assignment ranges (via interval arithmetic on all RHS expressions)
  intersected with its declared range.  The interval operations cover
  arithmetic (\texttt{range-+}, \texttt{range-*}), bitwise
  (\texttt{range-\&}, \texttt{range-<<}), and extraction
  (\texttt{bits}).

\item \textbf{Type proposal} (\texttt{propose-types!}).
  Calls \texttt{get-smallest-type} on each variable's converged range
  to choose the narrowest concrete type.  Unsigned is preferred;
  signed is used only when the range includes negative values.
  Variables exceeding 128~bits fall back to dynamic (unbounded).
\end{enumerate}

The code generator then classifies each integer into three categories:
\begin{description}
\item[\texttt{native}] --- fits in a machine word ($\leq 64$ bits).
  Assignments may need bit-masking (\texttt{Word.And(rhs, Mask)}).
\item[\texttt{wide}] --- fixed width $> 64$ bits.  Stored as
  \texttt{Mpz.T} with \texttt{ForceRange} calls.
\item[\texttt{dynamic}] --- unbounded \texttt{int}.  No truncation.
\end{description}

\subsubsection{Known limitations}

The current implementation has a documented limitation (noted in
\texttt{codegen-m3.scm}): it does not distinguish between
\emph{user-declared} widths (where truncation is intentional, per C
semantics) and \emph{compiler-inferred} widths (where truncation
would indicate a compiler bug).  Both cases use the same masking
logic.  In a verified framework, these two cases should be separated:
user-declared truncation is a design choice that the protocol checker
accepts, while inferred truncation that actually fires would be a
violation of the range analysis soundness.

Additionally, the interval arithmetic for bitwise operations is sound
but not always tight (conservative over-approximation), and there is
a noted potential bug in the handling of \texttt{bits} as an LHS
target (\texttt{interval.scm}).

\subsubsection{Type narrowing in the verified framework}

The verified framework does not replace \texttt{cspc}'s existing
range analysis---it wraps it with explicit proof obligations:

\begin{enumerate}
\item \textbf{Width consistency.}  Every assignment, arithmetic
  operation, and channel communication is type-correct at the
  proposed widths.  The range analysis already establishes this; the
  verification step makes the argument explicit and checkable.

\item \textbf{Protocol preservation.}  The protocol on each channel
  is preserved under the width change.  For simple
  \texttt{\{ v \}*} protocols this is automatic (the sequence of
  communications is unchanged, only the value domain narrows).

\item \textbf{Truncation audit.}  For each masking operation in the
  generated code, the tool reports whether it is (a)~provably
  unreachable (range analysis guarantees no overflow), (b)~intentional
  truncation (user declared a narrow type), or (c)~a potential problem
  (inferred type may be too narrow).
\end{enumerate}

\subsubsection{Structure refinement}

A common data refinement in the Fulcrum CSP toolchain is packing
multiple fields into a single wide channel.  The MiniMIPS Execute
stage sends a 107-bit structure:

\begin{Verbatim}[samepage=true]
structure EmData = (int(32) result; int(32) store_data;
                    int(32) addr; int(5) dest;
                    int(1) rw; int(1) mem_read;
                    int(1) mem_write; int(1) taken;
                    int(32) target);
\end{Verbatim}

The \texttt{cspc} compiler handles structure packing in compile3!
(\texttt{insert-pack-copies}, \texttt{insert-unpack-copies}), which
introduces temporaries for serialization.  These temporaries start as
unbounded \texttt{int} and are narrowed in compile6! using the
struct's total bit width (via \texttt{structdecl-width}).

The proof obligation is that the packing and unpacking are inverses:
if the sender packs fields into the structure and the receiver unpacks
them, the values are preserved.  The range analysis provides this
guarantee when the struct fields' ranges fit within their declared
widths.

%----------------------------------------------------------------------
\subsection{Stage 4: Implementation}
\label{sec:stage4}
%----------------------------------------------------------------------

The final stage maps the verified CSP design to a physical
implementation.  We describe three target paths.

\subsubsection{Bundled-data asynchronous circuits}
\label{sec:bd}

\textbf{Path:} \texttt{.csp} $\to$ \texttt{csp2verilog} $\to$
SystemVerilog with BD channel protocol.

Bundled-data circuits use a request-acknowledge handshake alongside
the data wires.  The timing assumption is that data is stable before
the request signal arrives---a \emph{matched delay} constraint that
must be satisfied by the physical implementation but is outside the
scope of CSP-level verification.

The \texttt{csp2verilog} tool generates SystemVerilog modules where
each CSP process becomes a finite state machine.  Channel operations
become handshake transactions:

\begin{center}
\begin{tabular}{ll}
\toprule
\textbf{CSP operation} & \textbf{BD protocol} \\
\midrule
\texttt{R!x} & Assert \texttt{req}, drive \texttt{data}; wait for
  \texttt{ack}; deassert \texttt{req}; wait for $\neg$\texttt{ack} \\
\texttt{L?x} & Wait for \texttt{req}; latch \texttt{data}; assert
  \texttt{ack}; wait for $\neg$\texttt{req}; deassert \texttt{ack} \\
\bottomrule
\end{tabular}
\end{center}

\textbf{Verified properties.}
The CSP-level protocol conformance guarantees that every handshake
completes: the sender always asserts \texttt{req} when the protocol
demands a send, and the receiver always asserts \texttt{ack} when the
protocol demands a receive.  What the CSP level does \emph{not}
guarantee is the timing: the matched-delay constraint is a physical
property that must be verified by static timing analysis.

\textbf{Slack assignment.}
The \texttt{slacker} tool computes the optimal slack (number of
pipeline buffers) for each channel to meet throughput targets.  In the
verified framework, each slack assignment is a pipeline insertion step
(Section~\ref{sec:refinement-steps}) with a trivial proof obligation.

\subsubsection{Synchronous RTL (dehandshaking)}
\label{sec:rtl}

\textbf{Path:} \texttt{.csp} $\to$ dehandshaking compiler $\to$
SystemVerilog with clock boundaries replacing handshakes.

Dehandshaking converts an asynchronous CSP design to synchronous RTL
by replacing each handshake with a clocked register transfer.  Each
channel becomes a registered pipeline stage; the handshake protocol
is replaced by valid/ready flow control (or simply pipeline registers
if the pipeline is fully balanced).

\textbf{Mapping.}
\begin{center}
\begin{tabular}{ll}
\toprule
\textbf{CSP concept} & \textbf{Synchronous RTL} \\
\midrule
Channel & Pipeline register (or valid/ready interface) \\
Process & Combinational logic between registers \\
Slack & Pipeline depth \\
Channel dependency & Data hazard \\
Guard (probe) & Stall signal \\
\bottomrule
\end{tabular}
\end{center}

\textbf{Verified properties.}
The CSP-level protocol conformance guarantees the correct sequencing
of data through the pipeline: if the CSP design is deadlock-free, the
synchronous implementation cannot stall indefinitely (assuming the
clock runs and resets complete).  The dehandshaking transformation is
a verified step: the tool must show that the RTL pipeline behavior
(under valid/ready flow control) is a refinement of the CSP channel
protocol.

\textbf{Advantages over direct RTL design.}
Starting from a verified CSP specification and dehandshaking gives the
designer the benefits of CSP-level reasoning (compositional
deadlock-freedom, protocol types) with a synchronous implementation
that can be processed by standard EDA tools.  Timing closure is
handled by the standard synchronous flow (synthesis, place-and-route,
STA).

\subsubsection{QDI pipeline circuits}
\label{sec:qdi}

\textbf{Path:} \texttt{.csp} $\to$ \texttt{csp2gma} $\to$
production rules $\to$ transistor-level implementation.

Quasi-delay-insensitive (QDI) circuits use delay-insensitive encoding
(typically dual-rail) and completion detection (C-elements) to
eliminate all timing assumptions except the isochronic fork.  This is
the most robust implementation style but also the most
area-intensive.

\textbf{Pipeline templates.}
Lines' pipelining theory~\cite{lines1998} defines a family of pipeline
templates with varying performance/area tradeoffs:

\begin{description}
\item[WCHB (Weak-Condition Half Buffer).]
  The simplest template.  Each stage is a half buffer that
  alternately receives and sends.  Throughput: 1 token per 2 cycles.
  Area: minimal.

\item[PCHB (Pre-Charge Half Buffer).]
  Uses pre-charged logic for higher speed.  Throughput: 1 token per
  2 cycles (faster cycle).  Area: moderate.

\item[PCFB (Pre-Charge Full Buffer).]
  A full buffer that can hold a token while simultaneously receiving
  the next.  Throughput: 1 token per 1 cycle (pipelined).
  Area: larger.
\end{description}

\textbf{Verified properties.}
Each pipeline template is a fixed process that implements a known
protocol.  The verification obligation is that the template's protocol
matches the channel protocol declared in the \texttt{.sys} file.
Since the templates are parameterized only by data width, this check
need only be done once per template type, not per instance.

The \texttt{csp2gma} tool generates a GMA intermediate
representation from which production rules can be derived.  The GMA
preserves the CSP-level semantics; the production rule derivation is a
separate verified step that must ensure the dual-rail encoding
faithfully represents the original data values.

%======================================================================
\section{Required Tools and Development Gaps}
\label{sec:tools}
%======================================================================

%----------------------------------------------------------------------
\subsection{Existing Tools (Ready to Use)}
\label{sec:existing}
%----------------------------------------------------------------------

The following tools are already implemented and available:

\begin{description}
\item[\texttt{cspc} / \texttt{cspfe} / \texttt{cspbuild}.]
  The core compilation and build pipeline.  Parses \texttt{.csp} and
  \texttt{.sys} files, compiles to Modula-3 state machines, generates
  simulator binaries.

\item[\texttt{csp2verilog} / \texttt{Cast2Verilog}.]
  BD synthesis path.  Generates synthesizable SystemVerilog from CSP
  specifications with bundled-data handshake protocols.

\item[\texttt{csp2gma}.]
  QDI synthesis path.  Generates GMA intermediate representation for
  production rule derivation.

\item[\texttt{slacker}.]
  Slack optimization.  Computes optimal buffer placements via linear
  programming.  Integrates with \texttt{Cast2Verilog} for
  post-synthesis buffer insertion.

\item[MDD-based deadlock checker.]
  Symbolic state-space exploration using Multi-valued Decision Diagrams
  with saturation algorithm.  Scales to 1024-process systems.
  Available as \texttt{check-deadlock-saturation!} in the Scheme
  environment.

\item[Explicit-state deadlock checker with POR.]
  BFS/DFS exploration with partial-order reduction and memoization.
  Provides counterexample traces.  Available as
  \texttt{check-deadlock!}.

\item[LTS extraction.]
  Per-process labelled transition system extraction from compiled
  process code.  Used by both the explicit-state and MDD-based
  checkers.
\end{description}

%----------------------------------------------------------------------
\subsection{Tools Needed (Development Gaps)}
\label{sec:gaps}
%----------------------------------------------------------------------

\begin{center}
\begin{tabular}{llll}
\toprule
\textbf{Tool} & \textbf{Purpose} & \textbf{Effort} & \textbf{Priority} \\
\midrule
Protocol type checker & Verify process conforms to & Medium & Critical \\
                      & channel protocols & & \\
Protocol parser       & Parse protocol annotations & Small & Critical \\
                      & in \texttt{.sys} files & & \\
Dependency graph      & Check acyclicity of channel & Small & Critical \\
analyzer              & dependency graph & & \\
Refinement proof      & Verify decomposition steps & Large & High \\
checker               & preserve behavior & & \\
Data refinement       & Verify type narrowing & Medium & High \\
checker               & preserves protocols & & \\
Pipeline insertion    & Verify Lines-style & Medium & High \\
verifier              & transformations & & \\
Dehandshaking         & CSP $\to$ synchronous RTL & Large & Medium \\
compiler              & & & \\
Protocol inference    & Infer protocols from & Large & Low \\
                      & existing code (migration) & & \\
\bottomrule
\end{tabular}
\end{center}

Each tool is described in detail below.

\subsubsection{Protocol type checker}

\textbf{Purpose.} Given a compiled process and its declared port
protocols, verify that the process conforms to the protocols
(Definition~\ref{def:conformance}).

\textbf{Approach.} Extract the per-process LTS (already implemented)
and check that every path through the LTS, projected onto each port,
is a word in the protocol's language.  Since protocol types are
regular, this reduces to checking that the product of the LTS with
the protocol automaton (for each port) has no error states.

\textbf{Effort: Medium.}  The LTS extraction infrastructure exists.
The main work is implementing the protocol automaton construction and
the product-state exploration with error detection.  Estimated: 2--4
weeks.

\subsubsection{Protocol parser}

\textbf{Purpose.} Extend the \texttt{.sys} file parser to recognize
\keyword{protocol} and \keyword{conforms} clauses.

\textbf{Approach.}  Add productions to the existing Scheme-based
\texttt{.sys} parser.  Protocol types are regular expressions, so the
parser is straightforward.

\textbf{Effort: Small.}  Estimated: 1--2 weeks, including integration
with the existing \texttt{validate-system!} procedure.

\subsubsection{Dependency graph analyzer}

\textbf{Purpose.} Construct the channel dependency graph from the
process code and check for cycles.

\textbf{Approach.} For each process, analyze the control flow to
determine which channel operations depend on which other channel
operations.  This is a straightforward dataflow analysis on the
per-process AST.  Cycle detection is a standard graph algorithm
(Tarjan's SCC or simple DFS).

\textbf{Effort: Small.}  Estimated: 1--2 weeks.  The main subtlety is
handling guards that probe multiple channels simultaneously.

\subsubsection{Refinement proof checker}

\textbf{Purpose.} Given a specification process $P$ and its
decomposition $P_1 \| P_2$, verify that the decomposition is correct
(the external behavior is preserved).

\textbf{Approach.} There are two levels of checking:
\begin{enumerate}
\item \textbf{Lightweight:} Check that $P_1$ and $P_2$ each conform
  to their port protocols, that internal channels have matching
  protocols, and that the dependency graph is acyclic.  This is
  sufficient for deadlock-freedom but does not guarantee trace
  equivalence.

\item \textbf{Full:} Construct the composed LTS of $P_1 \| P_2$,
  hide the internal channels, and check trace equivalence (or
  failures equivalence) with the original $P$.  This requires the
  model checking infrastructure but is applied only to the local
  decomposition, not the entire system.
\end{enumerate}

\textbf{Effort: Large.}  The lightweight check reuses the protocol
type checker and dependency graph analyzer.  The full check requires
LTS composition and equivalence checking, which is substantial new
infrastructure.  Estimated: 2--4 months for the full version.

\subsubsection{Data refinement checker}

\textbf{Purpose.} Verify that narrowing a channel's width from $W_1$
to $W_2 < W_1$ preserves the protocol.

\textbf{Approach.} Perform integer range analysis on the process code
to verify that no value exceeds the narrowed width.  For channels
with data constraints in their protocols, verify that the constraints
are still satisfiable.  This is a standard abstract interpretation
problem.

\textbf{Effort: Medium.}  Estimated: 3--6 weeks.  Range analysis for
fixed-width integers is well-understood; the main work is integrating
it with the Fulcrum CSP AST and type system.

\subsubsection{Pipeline insertion verifier}

\textbf{Purpose.} Verify that inserting a pipeline buffer (WCHB,
PCHB, PCFB, or simple FIFO) on a channel preserves the protocol.

\textbf{Approach.}  Each pipeline template has a fixed LTS.  The
verifier checks that the template's LTS, composed with the channel
protocol automaton, admits exactly the same external behavior (modulo
buffering delay).  Since the templates are fixed and small, this
check can be done once per template type and cached.

\textbf{Effort: Medium.}  Estimated: 2--4 weeks.  The main work is
formalizing the pipeline templates as LTSs and implementing the
equivalence check.

\subsubsection{Dehandshaking compiler}

\textbf{Purpose.} Compile a CSP design to synchronous RTL by
replacing handshakes with clocked register transfers.

\textbf{Approach.}  Map each channel to a pipeline register with
valid/ready flow control.  Map each process to combinational logic
with registered outputs.  Generate SystemVerilog with appropriate
clock-domain crossing logic for channels that span clock domains.

\textbf{Effort: Large.}  This is a new synthesis backend, comparable
in scope to \texttt{csp2verilog}.  Estimated: 3--6 months.  However,
it can leverage much of the existing \texttt{csp2verilog}
infrastructure (AST traversal, variable analysis, code emission).

\subsubsection{Protocol inference}

\textbf{Purpose.} Infer protocol types from existing \texttt{.csp}
code that lacks annotations.  This enables migration of existing
designs to the verified framework.

\textbf{Approach.} Analyze the process code to extract the sequence
of channel operations on each port.  Generalize the observed sequences
into a regular expression (the protocol type).  This is essentially
the problem of learning a regular language from positive examples.

\textbf{Effort: Large.}  Protocol inference is useful but not
essential---designers can annotate protocols manually.  The main
risk is that inferred protocols may be too permissive (accepting
behaviors the designer did not intend) or too restrictive (rejecting
valid refinements).  Estimated: 2--4 months.

%----------------------------------------------------------------------
\subsection{Integration Plan}
\label{sec:integration}
%----------------------------------------------------------------------

We recommend a phased implementation:

\subsubsection{Phase 1: Protocol types and conformance checker}

\textbf{Components:} Protocol parser, protocol type checker, dependency
graph analyzer.

\textbf{Outcome:} The designer can annotate \texttt{.sys} files with
protocol types and receive immediate feedback on conformance and
deadlock-freedom (for acyclic systems).

\textbf{Effort:} 1--2 months.

\textbf{Value:} This phase alone provides compositional deadlock-freedom
checking that scales to arbitrarily large systems without state-space
explosion.  For acyclic systems, it completely replaces the need for
model checking.

\subsubsection{Phase 2: Refinement proof checker}

\textbf{Components:} Lightweight refinement checker (protocol-based),
full refinement checker (LTS-based, for cyclic subgraphs).

\textbf{Outcome:} The designer can decompose processes and have each
decomposition step verified.  The model checking infrastructure
(MDD, POR) serves as the backend for checking cyclic subgraphs.

\textbf{Effort:} 2--4 months.

\textbf{Value:} Enables a verified design flow where each step is
checked locally.  The prior work on MDD-based deadlock checking is
reused as a backend, not discarded.

\subsubsection{Phase 3: Pipeline and slack verification}

\textbf{Components:} Pipeline insertion verifier, integration with
\texttt{slacker}.

\textbf{Outcome:} Slack insertion and pipeline optimization are
verified transformations.  The designer can trust that performance
optimization does not break correctness.

\textbf{Effort:} 1--2 months.

\textbf{Value:} Closes the gap between the CSP-level correctness proof
and the optimized implementation.

\subsubsection{Phase 4: Dehandshaking compiler}

\textbf{Components:} Dehandshaking compiler, data refinement checker.

\textbf{Outcome:} A new implementation target (synchronous RTL) with
verified correctness.  The data refinement checker ensures that type
narrowing for synthesis does not introduce errors.

\textbf{Effort:} 3--6 months.

\textbf{Value:} Opens the Fulcrum CSP toolchain to synchronous design
targets, broadening its applicability.

%======================================================================
\section{Worked Example}
\label{sec:example}
%======================================================================

This section walks through a complete design from specification to
implementation, illustrating every stage of the verified refinement
flow.  The example is a simplified 4-stage pipeline with feedback,
inspired by the MiniMIPS.

%----------------------------------------------------------------------
\subsection{Step 1: Write the Specification}
%----------------------------------------------------------------------

We begin with a monolithic specification of a simple
fetch-compute-writeback processor.  It reads instructions from an
input channel, computes a result, and writes back to a register file.

\begin{Verbatim}[samepage=true]
// spec.csp — monolithic processor specification
int(32) regs[8];
int(32) instr, rs, rt, result;
int(3)  rd_addr, rs_addr, rt_addr;
int(1)  wen;

*[
  IN?instr ;
  rs_addr = (instr >> 6) & 0x7 ;
  rt_addr = (instr >> 3) & 0x7 ;
  rd_addr = instr & 0x7 ;
  rs = regs[rs_addr] ;
  rt = regs[rt_addr] ;
  result = rs + rt ;          // simplified: always ADD
  wen = (instr >> 9) & 0x1 ;
  [ wen -> regs[rd_addr] = result
  [] ~wen -> skip
  ] ;
  OUT!result
]
\end{Verbatim}

The system description with protocol annotations:

\begin{Verbatim}[samepage=true]
system SimplePipe;
  var in_ch  : channel(32) protocol {instr }*;
  var out_ch : channel(32) protocol {result }*;

  process Spec = "spec.csp"
    port in  IN  : channel(32) conforms { instr }*
    port out OUT : channel(32) conforms { result }*;
begin
  var cpu : Spec(IN => in_ch, OUT => out_ch);
end.
\end{Verbatim}

\textbf{Tool output:}
\begin{Verbatim}[samepage=true]
$ check-protocols SimplePipe
Protocol conformance:
  Spec.IN:  PASS (conforms to { instr }*)
  Spec.OUT: PASS (conforms to { result }*)
Protocol matching:
  in_ch:  PASS (sender and receiver agree on { v }*)
  out_ch: PASS (sender and receiver agree on { v }*)
Dependency graph: acyclic (1 node, 0 edges)
Result: DEADLOCK-FREE (by compositional theorem)
\end{Verbatim}

%----------------------------------------------------------------------
\subsection{Step 2: Annotate Channel Protocols}
%----------------------------------------------------------------------

The specification has simple protocols (\texttt{\{ v \}*} on each
channel).  As we decompose, we will introduce more channels, each
needing its own protocol annotation.  For the 4-stage decomposition:

\begin{Verbatim}[samepage=true]
system SimplePipe4;
  // External channels
  var in_ch  : channel(32) protocol {instr }*;
  var out_ch : channel(32) protocol {result }*;

  // Internal pipeline channels
  var fd_instr : channel(32) protocol {instr }*;
  var dc_rs    : channel(32) protocol {rs_val }*;
  var dc_rt    : channel(32) protocol {rt_val }*;
  var dc_rd    : channel(3)  protocol {rd_addr }*;
  var dc_wen   : channel(1)  protocol {wen }*;
  var ex_result: channel(32) protocol {result }*;
  var ex_rd    : channel(3)  protocol {rd_addr }*;
  var ex_wen   : channel(1)  protocol {wen }*;

  // Feedback: Writeback -> Decode (register file)
  var wb_wen   : channel(1)  protocol {wen }*;
  var wb_rd    : channel(3)  protocol {rd_addr }*;
  var wb_val   : channel(32) protocol {val }*;
  // ...
\end{Verbatim}

Every channel in the decomposed system has an explicit protocol.
The tool verifies protocol matching at each channel endpoint.

%----------------------------------------------------------------------
\subsection{Step 3: Decompose into 4 Processes}
%----------------------------------------------------------------------

\textbf{Fetch stage:}

\begin{Verbatim}[samepage=true]
// fetch.csp
int(32) instr;
*[ IN?instr ; fd_instr!instr ]
\end{Verbatim}

\textbf{Decode stage:}

\begin{Verbatim}[samepage=true]
// decode.csp
int(32) instr, rs_val, rt_val;
int(3) rs_addr, rt_addr, rd_addr;
int(1) wen;
// (register file is a separate server process)
*[
  fd_instr?instr ;
  rs_addr = (instr >> 6) & 0x7 ;
  rt_addr = (instr >> 3) & 0x7 ;
  rd_addr = instr & 0x7 ;
  wen = (instr >> 9) & 0x1 ;
  rf_rs_req!rs_addr ; rf_rs_resp?rs_val ;
  rf_rt_req!rt_addr ; rf_rt_resp?rt_val ;
  dc_rs!rs_val ; dc_rt!rt_val ;
  dc_rd!rd_addr ; dc_wen!wen
]
\end{Verbatim}

\textbf{Execute stage:}

\begin{Verbatim}[samepage=true]
// execute.csp
int(32) rs, rt, result;
int(3) rd;
int(1) wen;
*[
  dc_rs?rs ; dc_rt?rt ; dc_rd?rd ; dc_wen?wen ;
  result = rs + rt ;
  ex_result!result ; ex_rd!rd ; ex_wen!wen
]
\end{Verbatim}

\textbf{Writeback stage:}

\begin{Verbatim}[samepage=true]
// writeback.csp
int(32) result;
int(3) rd;
int(1) wen;
*[
  ex_result?result ; ex_rd?rd ; ex_wen?wen ;
  wb_wen!wen ; wb_rd!rd ; wb_val!result ;
  OUT!result
]
\end{Verbatim}

\textbf{Register file server:}

\begin{Verbatim}[samepage=true]
// regfile.csp
int(32) regs[8];
int(32) val;
int(3) addr;
int(1) wen;
*[
  // Serve read requests, then accept writeback
  rf_rs_req?addr ; rf_rs_resp!regs[addr] ;
  rf_rt_req?addr ; rf_rt_resp!regs[addr] ;
  wb_wen?wen ; wb_rd?addr ; wb_val?val ;
  [ wen -> regs[addr] = val
  [] ~wen -> skip
  ]
]
\end{Verbatim}

\textbf{Tool output for decomposition check:}
\begin{Verbatim}[samepage=true]
$ check-protocols SimplePipe4
Protocol conformance:
  Fetch.IN:         PASS    Fetch.fd_instr:    PASS
  Decode.fd_instr:  PASS    Decode.dc_rs:      PASS
  Decode.dc_rt:     PASS    Decode.dc_rd:      PASS
  Decode.dc_wen:    PASS    Decode.rf_rs_req:  PASS
  Decode.rf_rs_resp:PASS    Decode.rf_rt_req:  PASS
  Decode.rf_rt_resp:PASS
  Execute.dc_rs:    PASS    Execute.dc_rt:     PASS
  Execute.dc_rd:    PASS    Execute.dc_wen:    PASS
  Execute.ex_result:PASS    Execute.ex_rd:     PASS
  Execute.ex_wen:   PASS
  Writeback.ex_*:   PASS    Writeback.wb_*:    PASS
  Writeback.OUT:    PASS
  RegFile.rf_*:     PASS    RegFile.wb_*:      PASS
Protocol matching: ALL PASS (14 channels checked)
Dependency graph:
  Cycle detected: Decode -> Execute -> Writeback -> RegFile -> Decode
  Annotation required: cycle broken by lockstep protocol
    (Decode always reads before Writeback writes)
  Localized check: running MDD checker on cycle subgraph...
  MDD result: DEADLOCK-FREE (4 processes, 12 channels)
Result: DEADLOCK-FREE
  (acyclic portion by compositional theorem,
   cyclic portion by MDD model checking)
\end{Verbatim}

%----------------------------------------------------------------------
\subsection{Step 4: Data Refinement}
%----------------------------------------------------------------------

The specification uses \texttt{int(32)} for all values.  For
synthesis, we narrow to minimal widths:

\begin{Verbatim}[samepage=true]
// Width refinement
var dc_rd    : channel(3)  protocol {rd_addr }*;  // was 32
var dc_wen   : channel(1)  protocol {wen }*;      // was 32
var ex_rd    : channel(3)  protocol {rd_addr }*;  // was 32
var ex_wen   : channel(1)  protocol {wen }*;      // was 32
var wb_wen   : channel(1)  protocol {wen }*;      // was 32
var wb_rd    : channel(3)  protocol {rd_addr }*;  // was 32
// Data channels remain 32-bit
\end{Verbatim}

\textbf{Tool output:}
\begin{Verbatim}[samepage=true]
$ check-data-refinement SimplePipe4
Width analysis:
  dc_rd:  3-bit sufficient (max value 7, from 3-bit mask)
  dc_wen: 1-bit sufficient (max value 1, from 1-bit mask)
  ex_rd:  3-bit sufficient (passthrough from dc_rd)
  ex_wen: 1-bit sufficient (passthrough from dc_wen)
  wb_wen: 1-bit sufficient (passthrough)
  wb_rd:  3-bit sufficient (passthrough)
Protocol preservation: ALL PASS
  (all protocols are { v }*, width narrowing preserves protocol)
Result: DATA REFINEMENT VALID
\end{Verbatim}

%----------------------------------------------------------------------
\subsection{Step 5: Add Pipeline Slack}
%----------------------------------------------------------------------

For throughput, we add pipeline buffers.  The \texttt{slacker} tool
computes the optimal assignment:

\begin{Verbatim}[samepage=true]
$ slacker SimplePipe4.sys
Slack assignment (minimizing buffer area):
  fd_instr: slack 1 (default, no change)
  dc_rs:    slack 1
  dc_rt:    slack 1
  dc_rd:    slack 2  // extra slack to balance pipeline
  dc_wen:   slack 2
  ex_result: slack 1
  ex_rd:    slack 1
  ex_wen:   slack 1
  wb_*:     slack 1
  rf_*:     slack 1
Total buffer area: 142 bit-slots
\end{Verbatim}

Each slack insertion is a verified transformation:

\begin{Verbatim}[samepage=true]
$ verify-slack-insertion SimplePipe4
Slack insertion verification:
  dc_rd (slack 1 -> 2): PASS
    Buffer conforms to { v }* on both ports
    Protocol preserved (output sequence = input sequence, delayed)
  dc_wen (slack 1 -> 2): PASS
  All other channels: slack 1 (no change, trivially valid)
Result: SLACK INSERTION VALID
\end{Verbatim}

%----------------------------------------------------------------------
\subsection{Step 6: Generate BD SystemVerilog}
%----------------------------------------------------------------------

Finally, we generate the bundled-data implementation:

\begin{Verbatim}[samepage=true]
$ csp2verilog SimplePipe4
Generating SystemVerilog:
  SimplePipe4_Fetch.sv     (2 states, 2 channels)
  SimplePipe4_Decode.sv    (6 states, 8 channels)
  SimplePipe4_Execute.sv   (2 states, 6 channels)
  SimplePipe4_Writeback.sv (2 states, 6 channels)
  SimplePipe4_RegFile.sv   (4 states, 6 channels)
  SimplePipe4_top.sv       (top-level with channel instantiation)
  SimplePipe4_tb.sv        (testbench)
\end{Verbatim}

\textbf{Verification summary for the complete flow:}

\begin{center}
\begin{tabular}{llll}
\toprule
\textbf{Step} & \textbf{Check} & \textbf{Method} & \textbf{Result} \\
\midrule
1. Specification & Protocol conformance & Local per-process & PASS \\
2. Protocol annotation & Duality & Syntactic & PASS \\
3. Decomposition & Conformance + acyclicity & Local + graph & PASS* \\
  & Cyclic subgraph & MDD model check & PASS \\
4. Data refinement & Width analysis & Abstract interp. & PASS \\
5. Slack insertion & Buffer conformance & Template check & PASS \\
6. BD synthesis & Handshake protocol & By construction & PASS \\
\bottomrule
\end{tabular}
\end{center}

\noindent *The cyclic portion (Decode--Execute--Writeback--RegFile)
requires localized model checking; the acyclic portion is verified by
the compositional theorem.

The total verification cost is dominated by the localized model
checking of the 4-process cycle.  For the acyclic portions, protocol
conformance checking is linear in the size of each process's LTS,
and dependency acyclicity is linear in the number of channels.  This
stands in contrast to global model checking, which must explore the
product of all process state spaces.

%======================================================================
\section{Related Work}
\label{sec:related}
%======================================================================

%----------------------------------------------------------------------
\subsection{Session Types}
%----------------------------------------------------------------------

Session types, introduced by Honda, Vasconcelos, and
Kubo~\cite{honda1998}, assign types to communication channels that
describe the sequence of messages exchanged.  The type system
guarantees progress (no deadlock) and session fidelity (no protocol
violation).  Yoshida and Carbone~\cite{yoshida2012} extended session
types to multiparty protocols.

Our protocol types are a specialization of binary session types to the
hardware domain: point-to-point channels with fixed-width data and
no delegation.  The key simplification is that our protocols are
regular (no recursive types beyond Kleene star), which makes
conformance checking decidable by finite automaton inclusion rather
than the more complex type-checking algorithms needed for general
session types.

Kobayashi's deadlock-freedom types~\cite{kobayashi2006} are
particularly relevant.  Kobayashi assigns ``usages'' to channels that
specify the order and multiplicity of operations, then verifies
acyclicity of the dependency relation.  Our
Theorem~\ref{thm:deadlock-free} is essentially the hardware
specialization of Kobayashi's result, with the channel dependency
graph replacing the implicit dependency relation in the type system.

Padovani~\cite{padovani2014} further refined Kobayashi's approach,
handling more general topologies including some cyclic ones.  Our
approach to cycles (localized model checking) is more pragmatic but
less elegant than Padovani's type-theoretic treatment.

%----------------------------------------------------------------------
\subsection{Event-B Refinement}
%----------------------------------------------------------------------

Event-B~\cite{abrial2010} is a formal method based on stepwise
refinement of abstract machines.  Each refinement step is accompanied
by proof obligations (invariant preservation, convergence, guard
strengthening) that are discharged by automated or interactive
provers.  The Rodin platform~\cite{rodin} provides tool support for
Event-B development.

Our refinement framework shares Event-B's philosophy of
correctness-by-construction through stepwise refinement, but differs
in several ways:
\begin{itemize}
\item Event-B's proof obligations are first-order logic formulas
  about state invariants; ours are protocol conformance checks
  (regular language inclusion).
\item Event-B uses general-purpose provers (Atelier~B, SMT solvers);
  our checks are domain-specific and decidable.
\item Event-B does not specifically target concurrent hardware;
  our framework is designed for the CSP channel model.
\end{itemize}

The tradeoff is generality vs.\ automation: Event-B can express
arbitrary system properties but requires manual proof guidance; our
framework handles only protocol-related properties but checks them
fully automatically.

%----------------------------------------------------------------------
\subsection{Rely-Guarantee Reasoning}
%----------------------------------------------------------------------

Jones' rely-guarantee reasoning~\cite{jones1983} verifies concurrent
programs by specifying, for each thread, the \emph{rely} condition
(what the environment guarantees) and the \emph{guarantee} condition
(what the thread provides).  Each thread is verified independently
against its rely and guarantee.

Protocol types serve an analogous role in our framework: a process's
port protocol is its guarantee (the communication pattern it promises),
and the matching protocol on the other endpoint is its rely (the
communication pattern it expects from the environment).  Protocol
conformance is the local verification step that ensures the guarantee
is met.

%----------------------------------------------------------------------
\subsection{Tangram and Handshake Circuits}
%----------------------------------------------------------------------

Van Berkel's Tangram compiler~\cite{vanberkel1993} takes a CSP-like
language and compiles it to handshake circuits in a
correct-by-construction fashion.  Each language construct maps to a
fixed handshake component (sequencer, parallel component, arbiter),
and the composition of correct components yields a correct circuit.

Tangram's approach is elegant but restrictive: the designer is
limited to the constructs supported by the language, and the
generated circuits are often larger than hand-designed ones.  Our
framework is less restrictive---the designer writes arbitrary CSP
processes and has them verified against protocols---but it requires
more tool support.  In essence, Tangram proves correctness by
limiting the design space; we prove correctness by checking the
design.

%----------------------------------------------------------------------
\subsection{Martin's Synthesis Methodology}
%----------------------------------------------------------------------

Alain Martin's synthesis methodology~\cite{martin1990} derives QDI
circuits from CSP specifications via a sequence of
semantics-preserving transformations: handshaking expansion,
production rule derivation, transistor sizing.  Each transformation
is justified by a semantic argument (typically, trace equivalence or
failure equivalence in the CSP theory).

Our framework is, in a sense, the \emph{formalization and
mechanization} of Martin's methodology.  Where Martin relies on the
designer's skill to ensure each transformation is correct, we propose
tool-checked proof obligations at each step.  The protocol type system
provides the formal scaffolding that Martin's methodology assumes
implicitly.

The key extension beyond Martin's work is the compositional
deadlock-freedom theorem (Theorem~\ref{thm:deadlock-free}), which
enables local verification of large systems without global
state-space exploration.  Martin's methodology does not explicitly
address compositionality for large systems---each design is verified
as a whole.

%----------------------------------------------------------------------
\subsection{Data-Driven Decomposition (Wong)}
%----------------------------------------------------------------------

Wong's data-driven decomposition (DDD)~\cite{wong2004} automates the
most labor-intensive step in Martin's methodology: decomposing a
sequential CHP program into a network of fine-grained concurrent
processes, each implementable as a single PCHB pipeline stage.  DDD
analyzes data dependencies in the sequential program and partitions it
into concurrent modules, preserving trace equivalence with the
original.  The method was validated on the instruction-fetch unit of
the Lutonium asynchronous 8051 microcontroller~\cite{martin2003}.

DDD operates at exactly the boundary addressed by our ``process
decomposition'' refinement step (Section~\ref{sec:refinement-steps}):
it takes a monolithic specification and produces a composed system
whose correctness relies on the same trace-equivalence argument.  The
key difference is that DDD is fully automatic (it \emph{computes} the
decomposition), while our framework verifies designer-written
decompositions.  The two approaches are complementary: DDD could serve
as an automated decomposition engine within the verified framework,
with the protocol type checker verifying the result.  Conversely, the
protocol type system could provide the correctness certificate that
DDD's decomposition algorithm currently establishes by construction.

%----------------------------------------------------------------------
\subsection{CHP Verification}
%----------------------------------------------------------------------

Myers~\cite{myers2001} and Manohar~\cite{manohar2001} have developed
verification techniques for CHP (Communicating Hardware Processes),
the language used by Martin's group.  These include:
\begin{itemize}
\item Deadlock and livelock detection via BDD-based state-space
  exploration.
\item Timing verification for QDI circuits using Boolean constraint
  propagation.
\item Protocol verification for handshake expansion.
\end{itemize}

Our approach differs in that we aim for compositional verification
(avoiding global state-space exploration for acyclic systems) and
integrate verification into the design flow rather than applying it
post hoc.  However, the CHP verification work provides valuable
algorithms and techniques that our implementation can draw on,
particularly for the QDI pipeline verification path.

%----------------------------------------------------------------------
\subsection{Lines Pipelining Theory}
%----------------------------------------------------------------------

Lines~\cite{lines1998} developed a theory of asynchronous pipelining
that characterizes pipeline templates (WCHB, PCHB, PCFB) by their
handshake protocols and performance properties.  Each template is a
CSP process with known behavior, and the theory guarantees that
composing correct templates yields a correct pipeline.

Our framework incorporates Lines' theory directly: each pipeline
template is a fixed process with a known protocol, and pipeline
insertion is a verified transformation (Section~\ref{sec:refinement-steps}).
The verification obligation for pipeline insertion---that the template's
protocol matches the channel protocol---is a formalization of Lines'
correctness argument.

The key contribution beyond Lines' work is the integration with the
protocol type system, which enables pipeline insertion to be verified
\emph{within} the same framework as all other design steps, rather
than as a separate correctness argument.

%======================================================================
\section{Conclusion}
\label{sec:conclusion}
%======================================================================

This report has presented a framework for verified refinement-based
design in the Fulcrum CSP toolchain.  The key ideas are:

\begin{enumerate}
\item \textbf{Protocol types on channels.}  Each channel carries an
  annotation specifying the permitted sequence of communications.
  Protocol types are regular expressions over send/receive actions,
  simple enough for fully automatic checking yet expressive enough to
  capture the communication patterns in real hardware designs.

\item \textbf{Compositional deadlock-freedom.}  If every process
  conforms to its port protocols, every channel has matching protocols at
  its endpoints, and the channel dependency graph is acyclic, then the
  system is deadlock-free by construction---no state-space exploration
  needed.  Cycles in the dependency graph are handled by localized
  model checking, reusing the existing MDD and explicit-state
  infrastructure.

\item \textbf{Local proof obligations.}  Each design step---process
  decomposition, data refinement, channel splitting, pipeline
  insertion---generates a proof obligation that can be checked locally,
  without composing the entire system.  This makes verification cost
  proportional to the complexity of the design step, not the size of
  the system.

\item \textbf{Three implementation targets.}  The verified design can
  be mapped to bundled-data async circuits (via \texttt{csp2verilog}),
  synchronous RTL (via dehandshaking), or QDI pipeline circuits (via
  \texttt{csp2gma}), with the CSP-level proof carrying across to each
  target.
\end{enumerate}

The model checking work from Phases~1--3 of the formal verification
project---LTS extraction, explicit-state checking with POR, MDD-based
saturation---is not discarded; it serves as the \emph{backend} for
verifying the cyclic portions of designs that the compositional theorem
cannot handle directly.  The shift is from using model checking as the
\emph{primary} verification method to using it as a \emph{fallback}
for the cases where compositional reasoning alone is insufficient.

The practical development path is clear:
\begin{enumerate}
\item Phase~1 (protocol types + conformance checker) provides
  immediate value and has modest implementation effort.
\item Phase~2 (refinement proof checker) enables the full verified
  design flow.
\item Phase~3 (pipeline/slack verification) closes the optimization
  gap.
\item Phase~4 (dehandshaking compiler) opens a new implementation
  target.
\end{enumerate}

The vision is a toolchain where the designer's knowledge of why a
design is correct is captured in the annotations, checked by the tool,
and preserved across every step from specification to silicon.
Post-hoc model checking reconstructs what the designer already knows;
verified refinement captures it at the source.

%======================================================================
% Bibliography
%======================================================================
\begin{thebibliography}{99}

\bibitem{hoare1985}
C.A.R.\ Hoare,
\textit{Communicating Sequential Processes},
Prentice-Hall, 1985.

\bibitem{fdr4}
T.\ Gibson-Robinson, P.\ Armstrong, A.\ Boulgakov, and A.W.\ Roscoe,
``FDR3: A Modern Refinement Checker for CSP,''
\textit{Proc.\ TACAS}, Springer LNCS 8413, pp.\ 187--201, 2014.

\bibitem{fv-report}
M.\ Nystr\"om and Claude (Anthropic),
``Formal Verification for the Fulcrum CSP Toolchain: Feasibility
Study, Project Plan, and Implementation Report,''
Intel Corporation, 2025.

\bibitem{martin1990}
A.J.\ Martin,
``Programming in VLSI: From Communicating Processes to
Delay-Insensitive Circuits,''
in C.A.R.\ Hoare, ed., \textit{Developments in Concurrency and
Communication}, UT Year of Programming Series, Addison-Wesley,
pp.\ 1--64, 1990.

\bibitem{abrial2010}
J.-R.\ Abrial,
\textit{Modeling in Event-B: System and Software Engineering},
Cambridge University Press, 2010.

\bibitem{honda1998}
K.\ Honda, V.T.\ Vasconcelos, and M.\ Kubo,
``Language Primitives and Type Discipline for Structured
Communication-Based Programming,''
\textit{Proc.\ ESOP}, Springer LNCS 1381, pp.\ 122--138, 1998.

\bibitem{kobayashi2006}
N.\ Kobayashi,
``A New Type System for Deadlock-Free Processes,''
\textit{Proc.\ CONCUR}, Springer LNCS 4137, pp.\ 233--247, 2006.

\bibitem{padovani2014}
L.\ Padovani,
``Deadlock and Lock Freedom in the Linear $\pi$-Calculus,''
\textit{Proc.\ CSL-LICS}, ACM, pp.\ 72:1--72:10, 2014.

\bibitem{jones1983}
C.B.\ Jones,
``Tentative Steps Toward a Development Method for Interfering
Programs,''
\textit{ACM TOPLAS}, vol.\ 5, no.\ 4, pp.\ 596--619, 1983.

\bibitem{vanberkel1993}
K.\ van Berkel,
\textit{Handshake Circuits: An Asynchronous Architecture for VLSI
Programming},
Cambridge University Press, 1993.

\bibitem{myers2001}
C.J.\ Myers,
\textit{Asynchronous Circuit Design},
Wiley, 2001.

\bibitem{manohar2001}
R.\ Manohar and A.J.\ Martin,
``Slack Elasticity in Concurrent Computing,''
\textit{Proc.\ Mathematics of Program Construction}, Springer LNCS
2386, pp.\ 272--285, 2001.

\bibitem{lines1998}
A.\ Lines,
``Pipelined Asynchronous Circuits,''
\textit{Caltech CS Technical Report CS-TR-95-21}, revised 1998.

\bibitem{peled1993}
D.\ Peled,
``All from One, One for All: On Model Checking Using Representatives,''
\textit{Proc.\ CAV}, Springer LNCS 697, pp.\ 409--423, 1993.

\bibitem{rodin}
J.-R.\ Abrial, M.\ Butler, S.\ Hallerstede, T.S.\ Hoang, F.\ Mehta,
and L.\ Voisin,
``Rodin: An Open Toolset for Modelling and Reasoning in Event-B,''
\textit{Int.\ Journal on Software Tools for Technology Transfer},
vol.\ 12, no.\ 6, pp.\ 447--466, 2010.

\bibitem{yoshida2012}
N.\ Yoshida and L.\ Gheri,
``Multiparty Session Types Meet Communicating Automata,''
\textit{Proc.\ ESOP}, Springer LNCS 7211, pp.\ 194--213, 2012.

\bibitem{nystrom2025csp}
M.\ Nystr\"om,
``The Fulcrum CSP Language Reference Manual,''
Intel Corporation, 2025.

\bibitem{cardi2011}
C.\ Cardi,
\textit{Slack Elasticity and Performance Analysis of Timed Asynchronous
Circuits},
Ph.D.\ thesis, University of Southern California, 2011.

\bibitem{beerel2010}
P.A.\ Beerel, R.O.\ Ozdag, and M.\ Ferretti,
\textit{A Designer's Guide to Asynchronous VLSI},
Cambridge University Press, 2010.

\bibitem{wong2004}
C.G.\ Wong,
\textit{High-Level Synthesis and Rapid Prototyping of Asynchronous
VLSI Systems},
Ph.D.\ thesis, California Institute of Technology, 2004.

\bibitem{martin2003}
A.J.\ Martin, M.\ Nystr\"om, and C.G.\ Wong,
``Three Generations of Asynchronous Microprocessors,''
\textit{IEEE Design \& Test of Computers}, vol.\ 20, no.\ 6,
pp.\ 9--17, 2003.

\bibitem{cortadella2002}
J.\ Cortadella, M.\ Kishinevsky, A.\ Kondratyev, L.\ Lavagno,
and A.\ Yakovlev,
\textit{Logic Synthesis of Asynchronous Controllers and Interfaces},
Springer, 2002.

\end{thebibliography}

\end{document}
